\chapter{Related Work} 
\label{ch:related-work}

% Motivation for p2p network monitoring
Both the Bitcoin network and the \gls{ln} operate as decentralized, peer-to-peer (\gls{p2p}) networks governed by open-source protocols.
A fundamental characteristic of them is the absence of a central authority. 
Instead, thousands of independent nodes continuously exchange data to maintain consensus and route payments. 
Obtaining a complete view of the network's topology and traffic is inherently difficult.

Despite these observability challenges, the transparent nature of both networks has generated interest across various areas. 
Academic researchers, commercial enterprises, mining pools, and government institutions all have interest in 
monitoring, capturing, and persisting this \gls{p2p} network data. 

\autoref{tab:motivation-p2p-monitoring} shows these diverse motivations, 
which can be grouped into two categories: \textit{technical} and \textit{intelligence}.
Researchers and protocol developers primarily focus on the \textit{technical} objectives 
e.g., mapping topologies, measuring propagation performance, and identifying vulnerabilities. 
In contrast, commercial entities and law enforcement agencies operate on \textit{intelligence} objectives e.g., 
using network observability to optimize routing services, monetize analytics, or conduct forensic financial tracing.

\begin{table}[htpb]
    \centering
    \caption{Motivations for monitoring Bitcoin and Lightning \gls{p2p} network traffic, grouped by entity and objective.}
    \label{tab:motivation-p2p-monitoring}
    \begin{tabularx}{\textwidth}{@{} p{2.5cm} l X @{}}
        \toprule
        \textbf{Entity} & \textbf{Primary Objective} & \textbf{Specific Motivations} \\
        \midrule
        Researchers \newline \& Developers
         & Vulnerability discovery & Observing and simulating network-level attacks. \\
         & Measuring performance & Analysing \gls{isp} centralization, propagation delays, and bandwidth overhead. \\
         & Topology mapping & Mapping network topology; analysing the distribution of nodes, clients, and channels. \\
        & Protocol improvement & Validating new protocol proposals against empirical network data. \\
        & Debugging & Detecting bugs in client implementations; monitoring node adoption of consensus upgrades. \\
        & Network health & Detecting transaction spam and potential network partitions. \\ 
        \addlinespace
       
        Miners & Competitive advantage & Monitoring transaction relay and prioritizing the rapid discovery of newly mined blocks. \\  
        \addlinespace
        
        Companies & Commercial Intelligence & Monetizing network analytics; optimizing routing algorithms and liquidity services. \\
        \addlinespace
        
        Governments & Law Enforcement & Sanctions compliance; identifying node operations within restricted jurisdictions. \\
         & Forensics & Tracing the origin of illicit funds (e.g., ransomware); de-anonymizing users by linking transaction broadcast IPs to public keys. \\
        \bottomrule
    \end{tabularx}
\end{table}


\section{Bitcoin Network Data}

% Academic
Academic literature about \gls{p2p} network traffic is extensive. 
In 2013 \citeauthor{decker_information_2013}  % \citeyear{} does not work?
laid the foundation with their paper \citetitle{decker_information_2013} ~\cite{decker_information_2013}
by measuring the propagation delay of newly found blocks and subsequently the number of forks occurring.

Nowadays, academic institutions like the Karlsruhe Institute of Technology (KIT) ~\cite{neudecker_2026_14627526} 
provide live analysis of the \gls{p2p} network data in their \term{bitnodes} project.
The analysis looks at different topics like the version of the \term{User agents} 
indicating the distribution of client versions on the Bitcoin mainnet, 
the geographic distribution of nodes by country, propagation delays and much more~\cite{kit_bitcoin_monitoring_web}. 


% Independent researchers
Due to Bitcoins free and open source ethos a variety of independent researchers work actively on Bitcoin. 
A subset of those focus on the \gls{p2p} network traffic.
The researcher \citeauthor{0xb10c_map_2026} maintains the project \citetitle{0xb10c_map_2026} showing 
various projects (including \tech{ln-history}) that focus on \gls{p2p} network monitoring~\cite{0xb10c_map_2026}.
Early Bitcoin Core contributor \citeauthor{corallo_21ninja}, shows his research, 
which focuses on DNS seed metrics and reachable nodes, on this website\footnote{\url{https://21.ninja}}.
Additionally, researcher \citeauthor{lopp_satoshi_info} maintains the \term{statoshi} project\footnote{\url{https://statoshi.info}}, 
providing extensive network visualizations.
To minimize propagation latency and reduce the risk of stale block races, miners can use 
high-speed overlay networks outside the traditional Bitcoin \gls{p2p} topology, 
such as the \gls{fibre}~\cite{corallo_fibre}.


% Companies and Governments 
Blockchain Intelligence Companies like \term{Chainalysis} provide tracing services 
which are demanded by law enforcement and compliance professionals. 
By using \term{attribution}, a process where information like a name or an ip-address gets 
linked to a (bitcoin-)address, they often own unique and sensitive information.
In 2025, \citeauthor{lubbertsen_ghost_2025} conducted a case study 
evaluating the accuracy of \term{Chainalysis} tracing services, 
demonstrating a false-positive rate of less than \qty{0.5}{\percent}~\cite{lubbertsen_ghost_2025}.

The need for this kind of observability was recognized early by governmental and regulatory institutions. 
In the early 2010s, reports from the \citeauthor{federalbureauofinvestigation_wayback_2014}~\cite{federalbureauofinvestigation_wayback_2014} 
and the \citeauthor{europeancentralbank_virtual_2012}~\cite{europeancentralbank_virtual_2012} 
highlighted the dual nature of decentralized networks. 
Both institutions warned that while Bitcoin's pseudonymous design introduces significant money-laundering risks, 
its permanently public broadcast data and blockchain simultaneously offer forensic tracing opportunities.



\section{Lightning Network Gossip Data}

The \gls{ln}, which builds upon the Bitcoin network, requires similar network observability. 
Since \gls{ln} nodes must independently route payments across channels, 
every participant must maintain an up-to-date view of the network topology. 
This state is continuously synchronized via \gls{p2p} gossip messages, 
the technical details of which are detailed in \autoref{ch:background}.

The analysis of this gossip data has been pioneered by several academic researchers. 
Notably, \citeauthor{decker_lnresearch_2020} and Prof. Schmid
published foundational studies~\cite{zabka_empirical_2022, zabka_short_2022}, 
releasing a dataset of over \num{35} million gossip messages spanning 
from the network's inception in 2018 to 2023~\cite{decker_lnresearch_2020}. 
This dataset served as the empirical basis for subsequent topological research, 
including quantitative analyses by \citeauthor{atmanaviciute_quantitative_2025}~\cite{atmanaviciute_quantitative_2025} 
and geolocation studies by \citeauthor{valko_geolocated_2025}~\cite{valko_geolocated_2025}. 
Unfortunately, this foundational dataset contains significant temporal gaps, 
and active collection by \citeauthor{decker_lnresearch_2020} has stopped. 
This lack of continuous and reliable data is the primary motivation 
for \tech{ln-history}: to establish an independent, open-source infrastructure dedicated 
to the persistent collection of \gls{ln} gossip. 

Parallel to academic efforts, protocol engineers actively monitor gossip 
to debug and optimize network performance. 
For instance, \citeauthor{harvey-buschel_jharveyb_2026} developed the \tech{Gossip-Observer} tool
to conduct 24-hour telemetry experiments, establishing simultaneous connections
with every public node to analyse message propagation~\cite{harvey-buschel_gossip_2025}. 
This research has been conducted at the Technical University of Berlin in 
\citeyear{gogge_routing_2022} by \citeauthor{gogge_routing_2022}~\cite{gogge_routing_2022},
who measured routing convergence delays within the \gls{ln}.
Similarly, \gls{ldk} maintainer \citeauthor{mattcorallo_bluematts_2026} 
has documented inefficiencies within the gossip protocol in the codebase~\cite{lightningdevkit_rustlightning_2025} 
and has published network snapshots utilizing the \term{rapid-gossip-sync-server}~\cite{corallo_lightningdevkit_2025}.

Finally, commercial interest in \gls{ln} data is driven by routing node operators, 
who require accurate routing metrics to allocate capital efficiently. 
Startups and network explorers such as \tech{mempool.space}, \tech{amboss.space}, 
and \tech{1ML.com} provide insights into the network topology. 
But to protect their business models, these entities often rely on closed-source 
algorithms and proprietary historical datasets. 

In contrast, \tech{ln-history} is engineered as a strictly open-source research platform. 
By providing a transparent foundation for historical and real-time data analysis, 
it empowers the community to drive protocol development by empirical data analysis.

